\chapter*{前言}
\addcontentsline{toc}{chapter}{前言}
\markboth{前言}{前言}
\SourcePage{6}
\begin{flushright}
\textit{一切都比人们想象的更简单,}\\
\textit{同时又比人们所能理解的复杂.}\\
Johann Wolfgang von Goethe\par\medskip
\textit{要到达未知之地,}\\
\textit{就得走一条自己不熟悉的路.}\\
St.~John of the Cross
\end{flushright}

本书讨论偏微分方程(PDE)的数值逼近.我们将详尽阐述各种数值方法,尤其是由偏微分方程的变分形式导出的方法,分析这些方法的稳定性和收敛性,推导误差界,并讨论实现这些方法所涉及的算法问题.

我们最重视的是在理论分析、算法介绍和应用讨论三者之间取得恰当的平衡.

本书涉及多类问题,包括线性与非线性问题、定常与含时问题,以及解光滑或不光滑的问题.除模型方程外,我们还考虑了若干在不同应用领域中具有实际意义的边值问题和初边值问题.

\BookRef{part:1}{第一部分}介绍并分析偏微分方程离散化的一般数值方法.

对于空间离散,我们将系统阐述 Galerkin 方法及其变体(Petrov--Galerkin 方法和广义 Galerkin 方法)的理论,以及配置法的理论.随后,我们将这一理论具体用于两类重要的近似子空间：一类来自有限元法,包括协调、非协调、混合和杂交有限元法；另一类来自谱方法,包括 Legendre 展开和 Chebyshev 展开.

对于非定常问题,我们将介绍用来进行时间推进的有限差分格式和分步格式.\BookRef{part:2}{第二部分}和\BookRef{part:3}{第三部分}还将在对流--扩散问题和双曲型方程的讨论中广泛涉及有限差分法.对于双曲型方程,我们也将简要介绍基于有限体积法的格式.

数值离散得到的代数方程组通常规模很大,而且非常稀疏.为求解这类方程组,我们回顾适用于对称和非对称矩阵的经典及现代技术,包括直接法与带预处理的迭代法.\SourcePage{7}我们还将简要介绍多重网格法和区域分解法.

\BookRef{part:2}{第二部分}和\BookRef{part:3}{第三部分}分别讨论定常问题和非定常问题.对于所考虑的每个边值问题或初边值问题,我们先介绍有关适定性的主要理论结果,即存在性、唯一性和先验估计方面的结果；随后回到前面介绍的数值方法,分析它们在当前问题中的表现,推导相应的代数形式,并讨论求解算法.

首先,我们考虑数学物理中的各类经典方程：描述势问题的椭圆型方程、描述热扩散的抛物型方程,以及描述波传播现象的双曲型方程.此外,我们还详细讨论被动标量的对流--扩散方程,以及描述黏性不可压缩流动的 Navier--Stokes 方程及其线性化形式(Stokes 问题).我们还将推导流体动力学方程的一般形式.

遗憾的是,受篇幅和我们自身经验所限,许多原本希望收入本书的重要课题未能纳入.例如,浅水方程的 Saint-Venant 模型,线弹性方程组及描述膜位移和薄板弯曲的双调和方程,半导体器件的漂移--扩散模型和流体动力学模型,以及可压缩流动的 Navier--Stokes 方程和 Euler 方程.

本书面向研究生,也面向从事偏微分方程数值模拟的研究人员和专业人员.

用作博士生课程教材时,可以讲授全书.\BookRef{part:1}{第一部分}可用作一个学季的偏微分方程变分数值方法入门课程.\BookRef{part:2}{第二部分}和\BookRef{part:3}{第三部分}分别将这些方法用于与时间无关和与时间有关的问题的数值逼近,可在余下两个学季讲授.其他安排也可行.例如,在\BookRef{part:1}{第一部分}的基础上加上第~\BookRef{ch:6}{6}~章、第~\BookRef{ch:11}{11}~章和第~\BookRef{ch:14}{14}~章的大部分内容,可组成一学期的课程,第二学期再讲授余下的内容.也可以按另一条主线组织教学,把\BookRef{part:1}{第一部分}同第~\BookRef{ch:8}{8}、\BookRef{ch:9}{9}、\BookRef{ch:10}{10}、\BookRef{ch:12}{12}、\BookRef{ch:13}{13}、\BookRef{ch:14}{14}~章合在一起,用作数值流体动力学导论.当然还可以作其他组合.

感谢 Springer-Verlag 的 C.~Byrne 博士和 J.~Heinze 博士在本书编写期间始终给予鼓励.Springer-Verlag 的技术人员也协助我们完成了书稿的定稿工作.

本书得益于我们过去几年在多所高等院校讲授这些课题的经历,以及学生的反馈.第一作者曾在位于明尼阿波利斯的明尼苏达大学、布雷西亚天主教大学和米兰理工大学讲授这些内容,第二作者则在特伦托大学讲授.\SourcePage{8}几位朋友和合作者也给予了帮助,他们或审阅了部分书稿,或提供了插图和表格.在此,我们向 V.~I.~Agoshkov、Yu.~A.~Kuznetsov、D.~Ambrosi、L.~Bergamaschi、S.~Delladio、M.~Manzini、M.~Paolini、F.~Pasquarelli、L.~Stolcis、E.~Zampieri、A.~Zaretti 致谢,尤其感谢 C.~Bernini、P.~Gervasio 和 F.~Saleri.

我们还要感谢 R.~Holliday 女士对整部书稿所作的语言编辑.最后,C.~Foglia 女士以精湛的技艺和令人惊叹的熟练程度录入了 \TeX{} 文件,这项工作对本书极为重要.

\noindent
\begin{tabular*}{\linewidth}{@{}l@{\extracolsep{\fill}}r@{}}
米兰和特伦托 & Alfio Quarteroni\\
1994 年 5 月 & Alberto Valli
\end{tabular*}

\bigskip
本书第二次印刷时,我们订正了若干印刷错误,并对原文作了一些修改.

具体来说,我们略微改动了第~\BookRef{sec:2.3.4}{2.3.4}、\BookRef{sec:3.4.1}{3.4.1}、\BookRef{sec:8.4}{8.4}、\BookRef{sec:12.3}{12.3}~节,并在笔记~\BookRef{note:8.2.1}{8.2.1}~中增补了一些说明.

我们还补全了 1994 年之后发表的论文的参考文献信息.

\noindent
\begin{tabular*}{\linewidth}{@{}l@{\extracolsep{\fill}}r@{}}
米兰和特伦托 & Alfio Quarteroni\\
1996 年 12 月 & Alberto Valli
\end{tabular*}
